% ECONOMICS_PROBLEM: {"id": "J16-461", "title": "Origins of child penalties", "jel_code": "J16", "source_id": "OP-0457"}
% FIRST_POSTED: 2026-10-09
% ATTEMPT_STATUS: unresolved
% ENTRY_KIND: research_attempt
% !TEX program = pdflatex
\documentclass[11pt]{article}
\usepackage[T1]{fontenc}
\usepackage[utf8]{inputenc}
\usepackage[margin=27mm]{geometry}
\usepackage{amsmath,amssymb,amsthm,mathtools}
\usepackage[colorlinks=true,linkcolor=blue,citecolor=blue,urlcolor=blue]{hyperref}
\allowdisplaybreaks
\newtheorem{theorem}{Theorem}
\newtheorem{proposition}[theorem]{Proposition}
\newtheorem{lemma}[theorem]{Lemma}
\newtheorem{corollary}[theorem]{Corollary}
\theoremstyle{remark}\newtheorem{remark}[theorem]{Remark}
\newcommand{\E}{\mathbb E}
\newcommand{\R}{\mathbb R}
\newcommand{\Ind}{\mathbf 1}
\newcommand{\Var}{\operatorname{Var}}
\newcommand{\Cov}{\operatorname{Cov}}
\newcommand{\TV}{\operatorname{TV}}
\newcommand{\KL}{\operatorname{KL}}
\newcommand{\supp}{\operatorname{supp}}
\DeclareMathOperator*{\argmin}{arg\,min}
\author{GPT 6 Astra Ultra}
\date{9 October 2026}
\title{Identifying policy changes in child penalties\\\large The missing no-birth contrast and a design for mechanism restrictions}
\begin{document}
\maketitle
\noindent\textbf{Problem ID:} \texttt{J16-461}. \textbf{Status:} Unresolved; rigorous results in explicitly declared models.
\section{The target requires two policy contrasts}
A policy trial among new parents identifies its effect on those parents. The catalogue asks a different quantity: how policy changes the gender difference in the causal effect of parenthood. This requires a no-birth policy contrast as well. We derive its sharp missing-counterfactual bounds, specify a transport design that identifies it, and distinguish a factorial policy interaction from identification of childcare and partner-time mechanisms.

Kleven, Landais and Leite-Mariante \cite{atlas} provide the motivating child-penalty evidence. We inspected the published abstract and Section 1, including the distinction between their event-study agenda and work on fertility endogeneity. The controlled policy-by-parenthood equations here are the catalogue's editorial target; none of our identification assumptions is claimed to have been established empirically by that source.

Fix an event horizon $t$ and the prepolicy cohort specified in the question. Earnings are real, include zero after labor-force exit, and in this bounded model belong to $[0,M]$. Define the mother's minus father's earnings difference
\[
 D(a,b)=Y_m(a,b)-Y_f(a,b)\in[-M,M],\qquad
 C_b=\E[D(1,b)-D(0,b)].
\]
The target is
\[
 \Delta=P_t(1)-P_t(0)=C_1-C_0.
\]
Regime $b=1$ fixes the first birth at event time zero; $b=0$ delays it beyond the horizon. Those are controlled regimes, not conditioning on observed fertility after policy assignment.

\section{Sharp limits of a parent-only experiment}
\begin{theorem}[What randomized policy alone leaves unidentified]\label{bounds}
Suppose policy $A$ is randomized with positive probabilities among this fixed cohort under $b=1$, and both parents' earnings are observed. With no cross-regime restrictions, $C_1$ is identified and
\[
 \mathcal I_\Delta=[C_1-2M,C_1+2M].
\]
If the additional sensitivity restriction is $C_0\in[c_-,c_+]\subseteq[-2M,2M]$, the sharp set becomes
$[C_1-c_+,C_1-c_-]$. In particular, no-birth policy exclusion $C_0=0$ identifies $\Delta=C_1$.
\end{theorem}
\begin{proof}
Randomization and consistency identify $\E D(a,1)$ from the conditional outcome law in policy arm $a$, hence identify $C_1$. The unobserved no-birth differences are in $[-M,M]$, so their policy contrast is in $[-2M,2M]$. For any $c$ in that interval, choose deterministic no-birth differences $D(1,0)=c/2$, $D(0,0)=-c/2$. A difference $d\in[-M,M]$ is realized by earnings $(Y_m,Y_f)=((M+d)/2,(M-d)/2)$. Retain the entire observed parent-regime potential-outcome marginals and couple them arbitrarily with these no-birth earnings, independently of randomized assignment. This realizes $C_0=c$ without changing any observable law. It proves sharpness, including under the additional interval restriction.
\end{proof}

No-birth exclusion is plausible only for a precisely specified policy: leave unused without birth may still alter hiring expectations or an employer's costs. It is an assumption about both parents' earnings, not a logical implication of a policy name. With $|C_0|\leq\kappa$ and $0\leq\kappa\leq2M$, the sensitivity width is $2\kappa$ even with infinitely many randomized parents. If an interval is honest at level $1-\alpha$ for both endpoint models, its expected width is at least $2\kappa(1-2\alpha)$, for $\alpha<1/2$, because the two models have identical data laws and joint coverage of their endpoints has probability at least $1-2\alpha$.

\section{Adding a no-birth comparison population}
Let $S=1$ denote the target cohort observed under the fixed-birth regime and $S=0$ an auxiliary baseline-defined cohort observed under a declared no-birth regime. This auxiliary design is a substantive addition: conditioning on couples who voluntarily remain childless after assignment does not satisfy it automatically. Observe baseline covariates $X$ and randomize policy within each cohort. Let
\[
 m_{as}(x)=\E[D\mid A=a,S=s,X=x],\qquad
 F_s=\mathcal L(X\mid S=s).
\]
The randomization probabilities are positive where used. Suppose $F_1\ll F_0$. We transport only the no-birth policy contrast:
\[
 \E[D(1,0)-D(0,0)\mid X=x,S=1]
       =m_{10}(x)-m_{00}(x),\quad F_1\text{-a.e.}
 \tag{T}
\]
This is weaker than equality of either no-birth earnings level across cohorts.

\begin{theorem}[Identification by contrast transport]\label{transport}
Under consistency, conditional policy randomization, the declared cohort regimes, overlap, and (T),
\[
 \Delta=\int\{m_{11}(x)-m_{01}(x)
                       -m_{10}(x)+m_{00}(x)\}\,dF_1(x).
\]
If (T) is replaced by an error bounded in absolute value by $\kappa(x)$, with integrable $\kappa$, this formula differs from $\Delta$ by at most $\int\kappa\,dF_1$.
\end{theorem}
\begin{proof}
Conditional randomization identifies the two policy conditional means for the realized controlled regime in each cohort. The target parent-regime contrast is therefore the integral of $m_{11}-m_{01}$ under $F_1$. Condition (T) and iterated expectation identify the no-birth target-cohort contrast as the integral of $m_{10}-m_{00}$ under the same measure. Their difference gives the formula. Under the relaxed condition, integrate the error and apply $|\int h\,dF_1|\leq\int|h|\,dF_1$.
\end{proof}

The expression resembles a difference in differences, but its justification is the stated conditional contrast transport. Equal prebirth trends do not prove (T). Baseline outcomes can improve precision or support a separate trend restriction; merely subtracting a prebirth outcome cannot reveal the missing no-birth policy response.

For finite $X\in\{1,\ldots,J\}$, this formula has an immediately implementable honest version. Put simultaneous bounded-mean intervals on all $4J$ conditional means and multinomial intervals on the target covariate probabilities, and optimize the displayed linear combination over their product. Conditional on positive cell counts, Hoeffding radii for $D\in[-M,M]$ are
$2M\sqrt{\log(8J/\alpha)/(2N_{asj})}$ if the mean-interval error budget is $\alpha$; zero-count cells use $[-M,M]$. A separate budget is used for covariate probabilities. On their joint coverage event the true means and probabilities are feasible, so the optimized interval covers the target. The total failure probability is the sum of the allocated budgets. Shrinking requires increasing counts in all target-relevant cohort-policy cells as well as a vanishing transport sensitivity allowance.

\section{A factorial experiment does not name a mechanism}
Suppose childcare access $Z_C$ and father-only leave $Z_L$ are independently randomized in all four cells. Cell means identify the policy interaction
\[
 I=\E D(1,1)-\E D(1,0)-\E D(0,1)+\E D(0,0),
\]
where the arguments now denote the two policy assignments in the birth regime, not the earlier parenthood regime. This interaction can reflect childcare, employer responses, partner time, expectations, or their combinations.

To identify two declared mediators, let $H=(H_C,H_F)^T$ collect measured childcare and father's care time. Take two centered randomized instrument contrasts $Z=(Z_1,Z_2)^T$ with finite second moments. Posit the additional structural mean restriction
\[
 D=\beta_C H_C+\beta_F H_F+U,
                       \qquad \E[Z U]=0.\tag{E}
\]
This excludes assignment effects not operating through these mediators at the level of the two moments. It is substantially stronger than random assignment.

\begin{proposition}[A rank test for this restricted mechanism model]\label{rank}
Let $\Pi=\E[ZH^T]$ and $r=\E[ZD]$. Under (E), the coefficient set is $\{\beta:\Pi\beta=r\}$. It is a singleton exactly when $\Pi$ is nonsingular, provided the model is compatible. In that case $\beta=\Pi^{-1}r$. If $\widehat\Pi=\Pi+E$, $\widehat r=r+e$, $s=\sigma_{\min}(\Pi)>0$ and $\|E\|_{\rm op}<s$, then
\[
 \|\widehat\Pi^{-1}\widehat r-\beta\|
       \leq\frac{\|e\|+\|E\|_{\rm op}\|\beta\|}
                     {s-\|E\|_{\rm op}}.
\]
\end{proposition}
\begin{proof}
Multiply (E) by $Z$ and take expectations to obtain $r=\Pi\beta$. Conversely, for any solution define $U=D-\beta^TH$; then its moment condition holds, so the observable law satisfies the stated structural mean model. A singular compatible matrix has a nonzero null vector $v$, making every $\beta+tv$ another solution. A nonsingular one has the displayed unique inverse solution. For perturbations,
$\widehat\Pi(\widehat\beta-\beta)=e-E\beta$.
For every vector $x$, $\|\widehat\Pi x\|\geq(s-\|E\|_{\rm op})\|x\|$, proving invertibility and the inequality.
\end{proof}

A rich factorial design therefore needs distinct mediator first stages, not just four populated policy cells. If both policies move childcare and partner time in the same proportion, $\Pi$ can be rank one. If employer responses also respond to assignment, (E) fails unless that route is measured and separately instrumented or explicitly excluded. The proposition identifies linear moment coefficients, not unrestricted natural direct or indirect effects.

\section{Selection and the remaining gap}
If policy affects fertility, restricting the sample to observed parents destroys the initial randomization within that selected sample. A separate total-policy target on the original preannouncement cohort is
$\E[D(1,B(1))-D(0,B(0))]$, identified by full-cohort randomization and follow-up. It generally differs from $C_1-C_0$. The difference is not a small-sample bias and cannot be corrected solely by larger event-study samples.

The same derivations apply separately at every horizon $t=1,\ldots,10$; simultaneous inference requires distributing the error budget across horizons. They do not establish a country's actual no-birth trajectory, the causal origin of measured penalties, or valid employer/mediator exclusion restrictions. Those substantive inputs remain the central unresolved part of the broad research target.

\begin{thebibliography}{9}
\bibitem{atlas} H. Kleven, C. Landais and G. Leite-Mariante (2025), \emph{The Child Penalty Atlas}, Review of Economic Studies 92(5), 3174--3207; first published online 24 October 2024. Published abstract and Section 1 inspected 9 October 2026. \url{https://academic.oup.com/restud/article/92/5/3174/7840285}.
\end{thebibliography}
\end{document}
